%Let $x \in \pi_{n} X$ be a homotopy element of a spectrum $X$ which has Adams filtration $k$. We let $\widetilde{x} \in (\pi_{n, n+k} \nu (X))$ denote any element of $\pi_{n,n+k} (\nu(X))$ with the property that $\tau^{k} \widetilde{x} = \nu(x)$. Such an $\widetilde{x}$ exists by \Cref{lemm:adams-fil1}.
%\todo{RB: These lifts aren't unique, so this introduces a weird ambiguity.}
Our main goal in this section is to prove \Cref{thm:mod8-main-thm}, which was a key input to our proof of \Cref{thm:mainboundary} in \Cref{sec:FinishingStolz}.

\begin{thm}\label{thm:mod8-main-thm}
    Let $F^s \pi_k (C(8)) \subseteq \pi_k (C(8))$ denote the elements of $\HFt$-Adams filtration at least $s$. Then for $k \geq 126$, the image of the map \[F^{\frac{1}{5} k + 15} \pi_k (C(8)) \to \pi_{k-1} (\Ss)\] is contained in the subgroup of $\pi_{k-1} (\Ss)$ generated by the image of $J$ and the $\mu$-family.

    %For $k \geq 143$, the natural map \[F^{\frac{1}{2} k - \frac{109}{5}} \pi_k (\Ss/8) \to F^{\frac{1}{5} k + 21} \pi_k (\Ss/8),\] where $F^s \pi_k (\Ss/8)$ denotes the elements of $H\F_2$-Adams filtration at least $s$, is an isomorphism.
\end{thm}
%and \[F^{\frac{1}{5}k + 21} \pi_k (\Ss/8) \to \pi_k (L_{K(1)} \Ss/8),\]

We will prove \Cref{thm:mod8-main-thm} by combining the banded genericity techonology of \Cref{sec:BandVan} with the main result of \Cref{sec:Y}. Before we explain further, let us fix some notation.

\begin{cnv}
    In this section, synthetic spectra will always be taken with respect to $\HFt$, and we will denote $\Sigma^{*,*}\nu_{\HFt}( \Ss_2^\wedge )$ by $\Ss_2^{*,*}$. Similarly, we will let $\Ss_2$ denote the $2$-complete sphere.
\end{cnv}

%% \begin{comment}
%% \begin{cnv}
%%     In this section, we will let $\left(\Ss^{i,j}\right) ^{\wedge} _2 = \Sigma^{i-j} \nu_{\HFt} (\left(\Ss^j \right)^{\wedge} _2)$.
%% \end{cnv}
%% \end{comment}

\begin{ntn}
By the calculations of \Cref{prop:syn-toda-range}, we see that that there are classes $\wt{2} \in \pi_{0,1} \Ss_2^{0,0}$, $\wt{\eta} \in \pi_{1,2} \Ss_2^{0,0}$ and $\wt{\nu} \in \pi_{3,4} \Ss_2^{0,0}$ which satisfy relations $\tau \wt{2} = \nu (2) = 2$, $\tau \wt{\eta} = \nu (\eta)$ and $\tau \wt{\nu} = \nu (\nu)$. Moreover, we let $\wt{2^n} = \wt{2}^n$.
\end{ntn}

\begin{lem} \label{fact:2-complete-fine}
    The natural map \[ [\Ss_2^{a,b}, \Ss_2 ^{0,0}] \to \pi_{a,b} (\Ss_2 ^{0,0})\] is an isomorphism for all $a,b$.
  Furthermore, for any $n \geq 0$ there is an equivalence
    \[\nu C(2^n) \simeq \mathrm{cof}(\Ss_2^{0,1} \xrightarrow{\tau^{n-1} \wt{2}^n} \Ss_2^{0,0}). \]
\end{lem}

\begin{proof}
  The first claim follows from \cite[Proposition 5.6]{Pstragowski}, which implies that $\Ss^{0,0} _2$ is the $\nu \HFt$-localization of $\Ss^{0,0}$. 

  %The second claim follows from the fact that $\Ss_2^0$ is $\HFt$-nilpotent complete and \Cref{lemm:easy-e-compton}.
  To prove the second claim, we note that the cofiber sequence $ \Ss_2^0 \to C(2^n) \to \Ss_2^1 $
    is short exact on $\HFt$-homology, so by \Cref{lemm:syn-cof} it induces a cofiber sequence $\Ss_2 ^{0,0} \to \nu C(2^n) \to \Ss_2 ^{1,1}$.
    Thus $\nu C(2^n)$ is the cofiber of a map $\Ss_2 ^{0,1} \to \Ss_2 ^{0,0}$ whose image under the functor $\tau^{-1}$ is $2^n$.
    The result therefore follows from the fact that $\pi_{0,*} (\Ss_2)$ is $\tau$-torsion free.
\end{proof}

\begin{ntn}
  For convenience, we will use the following notation:
  $$ C(\tau^a \wt{2}^b) \coloneqq \mathrm{Cof}(\Ss_2^{0,b-a} \xrightarrow{\tau^a \wt{2}^b} \Ss_2^{0,0}). $$
\end{ntn}

%% Before we begin the proof of \Cref{thm:mod8-main-thm}, we note down a fact and fix some notation.

%% \begin{fact}
%% \end{fact}

%% \begin{proof}
%%     This follows from \cite[Proposition 5.6]{Pstragowski}, which implies that $\nu (\Ss^{\wedge} _2)$ is the $\nu \HFt$-localization of $\Ss^{0,0}$.
%% \end{proof}

%% \begin{ntn}
%%     Given $x \in \pi_{a,b} (\Ss^{0,0})$, we let $C(x)^{\wedge} _2$ denote the cofiber of the self-map $\Sigma^{a,b} \nu (\Ss^{\wedge} _2) \to \nu (\Ss^{\wedge} _2)$ corresponding to $x$ under the fact above.
%% \end{ntn}


%% \begin{comment}
%% The uniqueness in the above defintion follows from the following lemma.

%% \begin{lem}
%%     For $k \leq 13$, the groups $\pi_{k,k+s} \left(\Ss^{0,0}\right) ^{\wedge} _2$ are $\tau$-torsion free.
%% \end{lem}

%% \begin{proof}
%%     Since the $\HFt$-Adams spectral sequence converges strongly to homotopy of $\left(\Ss^0 \right) ^{\wedge} _2$, it follows from \Cref{cor:strong-conv} that the $\nu_{\HFt} \HFt$-Adams spectral sequence converges strongly to the bigraded homotopy of $\left(\Ss^{0,0} \right) ^{\wedge} _2$. Since there are no differentials in the $\HFt$-Adams spectral sequence through total degree $13$, it follows from \Cref{thm:synth-ASS} and strong convergecne that there is no $\tau$-torsion in the relevant groups.
%% \end{proof}
%% \end{comment}
%
%The main step in the proof of \Cref{thm:mod8-main-thm} is to establish a $v_1$-banded vanishing line of slope $\frac{1}{5}$ for $C(\wt{8})$ with explicit parameters. By \Cref{thm:Y-v1-band}, $\nu (Y) \simeq C(\wt{2}) \otimes C(\wt{\eta})$ admits a $v_1$-banded vanishing line of slope $\frac{1}{5}$. To prove that $C(\wt{8})$ does too, as well as find explicit parameters, we will apply \Cref{prop:band-in-cofiber-seqs} to an explicit sequence of cofiber sequences which exhibits $C(\wt{8})$ as living in the thick subcategory of synthetic spectra generated by $\nu (Y)$. Once we have done that, we will convert our banded vanishing line into \Cref{thm:mod8-main-thm} by invoking \Cref{thm:DM-filt-bound}.

\begin{rmk}\label{rmk:modified-mod8}
    The synthetic spectrum $C(\wt{2}^n)$ encodes the modified $\HFt$-Adams spectral sequence for $C(2^n)$. See \cite[Section 3]{BHHM} for the notion of a modified Adams spectral sequence. \todo{Is there a better reference to cite here?}
\end{rmk}

Our next goal will be to establish a $v_1$-banded vanishing line of slope $\frac{1}{5}$ for $C(\wt{8})$ with explicit parameters. We will do this via a thick subcategory argument.

%Before we establish the desired $v_1$-banded vanishing line for $C(\wt{8})$, we prove a technical lemma.

\begin{lem}\label{lemm:eta-splitting}
    There is a splitting of synthetic spectra
    \[C(\wt{2}) \otimes C(\wt{\eta}^3) \simeq C(\wt{2}) \oplus \Sigma^{4,6} C(\wt{2}).\]
\end{lem}

\begin{proof}
  After inverting $\tau$ this splitting becomes the classical fact that
  $ C(2) \otimes C(\eta^3) \simeq C(2) \oplus \Sigma^{4} C(2) $
  and the proof we give simply lifts this argument to the synthetic setting.
  
  The splitting follows two statements:
  that $\wt{\eta}^3 = \wt{4} \wt{\nu}$ as self-maps of $C(\wt{2})$ and
  that $\wt{4}$ is null as a self-map of $C(\wt{2})$.
  The first fact follows from \Cref{prop:syn-toda-range}, which shows that the relation $\wt{\eta}^3 = \wt{4} \wt{\nu}$ holds in the homotopy of $\Ss^{0,0} _2$.
  To prove that $\wt{4}$ is null as a self-map of $C(\wt{2})$, we examine the following commutative diagram
  
  \begin{center}
    \begin{tikzcd}
      \Ss^{0,0}_2 \arrow[r] \arrow[d, "\wt{2}"] & C(\wt{2}) \arrow[r] \arrow[d, "\wt{2}"] \arrow[dl, dashrightarrow] & \Ss^{1,1}_2 \arrow[d, "\wt{2}"] \\
      \Ss^{0,0}_2 \arrow[r] \arrow[d, "\wt{2}"] & C(\wt{2}) \arrow[r] \arrow[d, "\wt{2}"] & \Ss^{1,1}_2 \arrow[dl, dashrightarrow] \arrow[d, "\wt{2}"] \\
      \Ss^{0,0}_2 \arrow[r] & C(\wt{2}) \arrow[r] & \Ss^{1,1}_2 ,
    \end{tikzcd}
  \end{center}
  where the rows are cofiber sequences and the dashed arrows exist because of the canonical nullhomotopies of \[C(\wt{2}) \to \Ss^{1,1}_2 \xrightarrow{\wt{2}} \Ss^{1,1}_2 \text{ and } \Ss^{0,0}_2 \xrightarrow{\wt{2}} \Ss^{0,0}_2 \to C(\wt{2}).\]
    We wish to show that the composite of the middle two vertical arrows is null. Using the dashed arrows, we may factor this through the composition of the middle two horizontal arrows, which is null because they form a cofiber sequence.
\end{proof}

\begin{prop}\label{prop:banded-numbers}
  The synthetic spectra $X$ in the table below admit $v_1$-banded vanishing lines of slope $\frac{1}{5}$ and remaining parameters as follows:
  \renewcommand{\arraystretch}{1.6}
  \begin{center}
    \begin{tabular}{|c||c|c|c|c|c|}
      \hline
      \textup{Synthetic Spectrum X} & $b$ & $d$ & $v$ & $c$ & $r$ \\ \hline \hline
      $C(\wt{2}) \otimes C(\wt{\eta})$ &
      $-1.5$ & $0$ & $15$ & $2.6$ & $1$ \\ \hline
      $C(\wt{2}) \otimes C(\wt{\eta}^2)$ &
      $-2.5$ & $0.5$ & $23$ & $4.4$ & $2$ \\ \hline
      $C(\wt{2}) \otimes C(\wt{\eta}^3)$ &
      $-3.5$ & $1$ & $32 + \frac{1}{3}$ & $6.2$ & $4$ \\ \hline
      $C(\wt{2})$ &
      $-3.5$ & $1$ & $28 + \frac{1}{3}$ & $5$ & $4$ \\ \hline
      $C(\wt{4})$ &
      $-7.5$ & $2$ & $58 + \frac{1}{3}$ & $10$ & $9$ \\ \hline
      $C(\wt{8)}$ &
      $-12.5$ & $3$ & $91 + \frac{2}{3}$ & $15$ & $15$ \\ \hline
    \end{tabular}
  \end{center}
  \renewcommand{\arraystretch}{1.0}
\end{prop}

\begin{proof}
  Inductively apply \Cref{lemm:band-retract-susp} and \Cref{prop:band-in-cofiber-seqs} to the Bockstein cofiber sequences
  \begin{align*}
      \Sigma^{1,2} C(\wt{2}) \otimes C(\wt{\eta}) \to C(\wt{2}) &\otimes C(\wt{\eta}^2) \to C(\wt{2}) \otimes C(\wt{\eta}) \\
      \Sigma^{1,2} C(\wt{2}) \otimes C(\wt{\eta}^2) \to C(\wt{2}) &\otimes C(\wt{\eta}^3) \to C(\wt{2}) \otimes C(\wt{\eta}) \\
      \Sigma^{0,1} C(\wt{2}) \to &C(\wt{4}) \to C(\wt{2}) \\
      \Sigma^{0,1} C(\wt{4}) \to &C(\wt{8}) \to C(\wt{2})
  \end{align*}
    using \Cref{thm:Y-v1-band} as a base case and \Cref{lemm:eta-splitting} to go from the second to the third sequence.
\end{proof}

\begin{rmk}
    The numbers in \Cref{prop:banded-numbers} can likely be improved by more carefully accounting for the behavior of the classes in the band under the cofiber sequences used in the proof. In particular, we believe that one could improve the torsion order bound in the $v_1$-banded vanishing line for $C(\wt{2})$ to $3$. This would imply by \Cref{prop:banded-adams} that the $v_1$-localized Adams spectral sequence for $C(2)$ collapses at the $\mathrm{E}_5$-page. This result was announced by Mahowald \cite[Theorem 5]{MahBull}, but to the best of our knowledge a proof has never appeared in the literature.
\end{rmk}

\begin{prop} \label{prop:I-want-to-be-done}
  If $C(\wt{2}^n)$ admits a banded vanishing line with parameters $(m, c, r, b, d, v)$, then
  the image of the map
  $$ F^{mk+c}\pi_k C(2^n) \to \pi_{k-1}(\Ss) $$
  is contained in the subgroup of $\pi_{k-1}(\Ss)$ generated by the image of $J$ and the $\mu$-family as long as $ k \geq v$ and
  $$ \frac{1}{2}k + b - n + 1 \geq \frac{3}{10}(k-1) + 4 + v_2(k+1) + v_2(k). $$
  Recall that $v_2 (k)$ denotes the $2$-adic valuation of $k$.
\end{prop}

\begin{proof}
  %\Cref{prop:banded-numbers} implies that the desired results are true for the filtration on $\Ss/8$ induced by $S^{0,0} / \wt{8}$. It therefore suffices to prove that there is a map $\nu(\Ss/8) \to S^{0,0} / \wt{8}$, since filtrations may only rise under such a map.
  First, we note that the conclusion holds trivially for $k \leq 1$.
  Next, using that $\pi_{k-1}(\Ss) \cong \pi_{k-1}(\Ss_2)$ for $k > 1$ and that the $\HFt$-Adams filtrations on each group agree, we may replace $\Ss$ in the theorem statement by $\Ss_2$. 
%  Since the homotopy of $C(2^n)$ is $4^n$-torsion, the homotopy of $\Ss$ is equal to the homotopy of $\nu \Ss_2^\wedge$, we may as well 

  Consider the diagram below, where each row is a cofiber sequence and the middle and right vertical maps are projection onto the top cell:
  \begin{center}
    \begin{tikzcd}
      \Sigma^{0,n}C\tau^{n-1} \ar[r] \ar[d, equal] &
      C(\tau^{n-1} \wt{2}^n) \ar[r] \ar[d] &
      C(\wt{2}^n) \ar[d] \\
      \Sigma^{0,n}C\tau^n \ar[r] &
      \Ss_2^{1, 1} \ar[r, "\tau^n"] &
      \Ss_2^{1, n} 
    \end{tikzcd}
  \end{center}
  By \Cref{fact:2-complete-fine}, there is an equivalence $ \nu C(2^n) \simeq C(\tau^{n-1} \wt{2}^n) $.
  Note that under $\tau^{-1}$ the map $\nu C(2^n) \to \Ss^{1,n}$ becomes projection onto the top cell.
  This implies that projection to the top cell induces maps
  $$ F^{s} \pi_k( C(2^n)) \to F^s \pi_k (\tau^{-1} C(\wt{2}^n)) \to F^{s-n+1} \pi_{k-1} \Ss, $$
  where $ F^s \pi_k(\tau^{-1} C(\wt{2}^n)) $ is as in \Cref{dfn:synth-filt}.\footnote{ \Cref{rmk:modified-mod8} allows us to identify this filtration with the \emph{modified} $\HFt$-Adams filtration.}
  We finish by using the hypothesis that $C(\wt{2}^n)$ has a banded vanishing line.
  It follows that, for $k \geq v$, the maps induced by projection to the top cell factor as
  \begin{align*}
    F^{mk + c} \pi_k (C(2^n))
    &\to F^{mk + c} \pi_k (\tau^{-1} C(\wt{2}^n) )
    = F^{\frac{1}{2}k + b} \pi_k (\tau^{-1} C(\wt{2}^n) )\\
    &\to F^{\frac{1}{2}k + b - n + 1} \pi_{k-1} \Ss_2.
  \end{align*}
  
  It therefore suffices to find a $k \geq v$ large enough so that every element of $\pi_{k-1} \Ss$ which has $\HFt$-Adams filtration at least $\frac{1}{2}k + b - n + 1$ is in the subgroup generated by the image of $J$ and the $\mu$-family.
  
  \Cref{thm:gamma-upper}(1) states that every element in $\pi_{k-1} \Ss$ which has $\HFt$-Adams filtration at least $\frac{3}{10}(k-1)+4+v_2 (k+1) + v_2 (k)$ is in the subgroup generated by the image of $J$ and the $\mu$-family. The result follows.
%  It therefore follows that it suffices to prove that \[\frac{1}{2}(k-1) - 14.5 \geq \frac{3}{10}(k-1)+4+v_2 (k+1) + v_2 (k)\] for all $k \geq 129$, which we do in the lemma below.
\end{proof}

\begin{proof}[Proof of \Cref{thm:mod8-main-thm}]
  Using Propositions \ref{prop:banded-numbers} and \ref{prop:I-want-to-be-done}, it will suffice to show that the following inequality holds for all $k \geq 126$:
  $$ \frac{1}{2}k - 14.5 \geq \frac{3}{10} (k-1) + 4 + v_2 (k+1) + v_2 (k). $$
  Rearranging, clearing denominators and applying the bound
  $$ \log_2(k+1) \geq v_2(k+1)+v_2(k), $$
  we find that it suffices to show that
  $$ k \geq 91 + 5\log_2(k+1). $$
  Taking derivatives, we find that the left hand side increases faster than the right hand side as soon as $k \geq 9$.
  Thus, to show the inequality holds for $k \geq 126$ it suffices to note that
  \[ 126 \geq 91 + 5 \log_2 (127) \approx 125.94. \qedhere\]

  % the desired inequality into \[\frac{1}{5}k \geq 18.3 + v_2 (k+1) + v_2 (k).\] Using the inequality $\log_2 (k) \geq v_2 (k)$, we find that it suffices to show that \[\frac{1}{5}k \geq 18.3 + 2 \log_2 (k+1),\] i.e. that \[k \geq 91.5 + 10 \log_2 (k+1).\]  Thus, to show the inequality holds for $k \geq 166$ it suffices to note that \[166 \geq 91.5 + 10 \log_2 (167) \approx 165.3.\]

  % Finally, we may use a computer to verify that the original inequality holds for $129 \leq k \leq 165$.
\end{proof}


%We conclude by proving a technical result used earlier in the paper:

%\begin{prop} \label{prop:trivial-square}
%  Fix an integer $n \geq 3$, and let $j \in \pi_{4n-1} \bS$ denote a generator of $\mathcal{J}_{4n-1}$.  Then the square of $j$ is trivial in $\pi_{8n-2} \bS$.
%\end{prop}

%\begin{proof}
  %% get j with filtration
%  In order to show that $j^2 = 0$ it will suffice to show that $j^2$ is in the image of $J$, because the image of $J$ is empty in stems of the form $8n-2$.
	
%  From \cite[Corollary 1.3]{DM3}, we know that $j$ has Adams filtration at least $2n - v_2(4n)$.
  %% sqaure rho and lose a filtration
%  Using \Cref{lemm:adams-fil1}, we may lift $j$ to a map
%  $$ \tilde{j} : \Ss^{4n-1, 6n-1-v_2(4n)} \to \Ss^{0,0}. $$
%	Using the fact that $\tilde{j}$ lives in odd topological degree, we learn from the $\mathbb{E}_\infty$-ring structure on $\Ss^{0,0}$ that $2\tilde{j}^2 = 0$.
%  Using the relation $2 = \tau \wt{2}$ we can then build a map
%  $$\Ss^{8n-1, 12n - 2 - 2v_2(4n)} \stackrel{f}{\longrightarrow} C(\wt{2}) \simeq \nu C(2) $$
%  such that composing $f$ with projection to the top cell gives $\tau \tilde{j}^2$.
%    Then \Cref{cor:synth-filt}\footnote{Note that this corollary applies because the $\HFt$-Adams spectral sequence converges strongly for the $2$-complete finite spectrum $C(2)$ by \cite[Theorem 6.6]{BousfieldLocalization}.} implies that image of $f$ under $\tau^{-1}$ has Adams filtration at least $4n - 2v_2(4n) - 1$. In particular, this means we can use Propositions \ref{prop:banded-numbers} and \ref{prop:I-want-to-be-done} to conclude that $j^2$ is in the image of $J$ as long as 
%  \begin{enumerate}
%  \item $ 8n-1 \geq 29 $,
%  \item $ 4n - 2v_2(4n) - 1 \geq \frac{1}{5}(8n-1) + 5 $, \text{ and}
%  \item $ \frac{1}{2}(8n-1) - \frac{7}{2} \geq \frac{3}{10}(8n-2) + 4 + v_2(8n) + v_2(8n-1). $
%  \end{enumerate}  
  %% inequality
%  After simplifying it suffices to show that
%  \begin{enumerate}
%  \item[($1'$)] $ n \geq 3.75 $,
%  \item[($2'$)] $ \frac{12}{5} n \geq 2\log_2(n) + 9.8 $, \text{ and}
%  \item[($3'$)] $ \frac{8}{5} n \geq \log_2(n) + 10.4 $.
%  \end{enumerate}
%  Taking derivatives, we find that, in both ($2'$) and ($3'$), the left hand side increases faster than the right hand side as soon as $n \geq 9$. Thus, in order to conclude for $n \geq 9$, it suffices to note that ($2'$) holds for $n = 7$ and ($3'$) holds for $n = 9$.
  
  %% remaining cases
%  For the remaining cases $n=3,4,5,6,7,8$, we note that the Adams filtration of $j^2$ is sufficiently high that we may conclude that $j^2 = 0$ from the low-%dimensional calculations of \cite{Isaksen}.
%\end{proof}


%%% Local Variables:
%%% mode: latex
%%% TeX-master: "Main"
%%% eval: (visual-line-mode t)
%%% eval: (auto-fill-mode 0)
%%% End:
